% =============================================================================
%  Page layout, headings, captions, code style
%  Check the Thesis C course outline for required margins / spacing / page limit
% =============================================================================
\usepackage[a4paper,top=25mm,bottom=25mm,left=30mm,right=25mm,headheight=15pt,marginparwidth=22mm]{geometry}
\onehalfspacing
\setlength{\parskip}{0.5em}
\setlength{\parindent}{0pt}

% ---- Headers / footers -------------------------------------------------------
\usepackage{fancyhdr}
\pagestyle{fancy}
\fancyhf{}
\fancyhead[L]{\small\nouppercase{\leftmark}}
\fancyfoot[C]{\thepage}
\renewcommand{\headrulewidth}{0.4pt}
\fancypagestyle{plain}{\fancyhf{}\fancyfoot[C]{\thepage}\renewcommand{\headrulewidth}{0pt}}

% ---- Chapter headings: compact, numbered ---------------------------------------
\usepackage{titlesec}
\titleformat{\chapter}[display]
  {\normalfont\bfseries\Huge}{\large\MakeUppercase{\chaptertitlename}~\thechapter}{0.5em}{}
\titlespacing*{\chapter}{0pt}{0pt}{1.5em}
\setcounter{secnumdepth}{3}
\setcounter{tocdepth}{2}

% ---- Lists --------------------------------------------------------------------
\setlist{itemsep=0.2em,topsep=0.3em}

% ---- Floats -------------------------------------------------------------------
\graphicspath{{figures/}}
\renewcommand{\floatpagefraction}{0.8}
\renewcommand{\topfraction}{0.9}
\renewcommand{\textfraction}{0.1}

% ---- Code listings (for short URScript / Python excerpts in appendices) -------
\lstdefinestyle{thesis}{
  basicstyle=\ttfamily\footnotesize, keywordstyle=\color{RoyalBlue}\bfseries,
  commentstyle=\color{ForestGreen}, stringstyle=\color{BrickRed},
  numbers=left, numberstyle=\tiny\color{gray}, numbersep=6pt,
  frame=single, rulecolor=\color{gray!40}, backgroundcolor=\color{gray!5},
  breaklines=true, showstringspaces=false, tabsize=4, captionpos=b}
\lstset{style=thesis}

% ---- Theorem-like environments (definitions of metrics, design laws) -----------
\theoremstyle{definition}
\newtheorem{definition}{Definition}[chapter]
\newtheorem{hypothesis}{Hypothesis}
\newtheorem{contribution}{Contribution}
